\documentclass[11pt,a4paper]{article}
\usepackage[margin=2.5cm]{geometry}
\usepackage{amsmath,amssymb}
\usepackage{booktabs}
\usepackage{parskip}
\usepackage{hyperref}
\usepackage{xcolor}
\usepackage{graphicx}

\hypersetup{colorlinks=true,linkcolor=blue,citecolor=blue}

\title{\textbf{Bootstrap Calibration of the Destacking Pipeline}\\[0.4em]
\large Calibration study --- destacking project}
\author{destacking project}
\date{2026 May 24}

\begin{document}
\maketitle
\tableofcontents
\newpage

% ============================================================
\section{Overview}
% ============================================================

The destacking pipeline constrains a magnitude-gap slope $\gamma$ by jointly
fitting $N = 5000$ per-cluster weak-lensing excess surface density ($\Delta\Sigma$)
profiles measured from simulated weak-lensing data.  Individual posterior
distributions on $\gamma$ are known to be \emph{under-dispersed}: they correctly
capture shape-noise (measurement) uncertainty but miss cluster sample variance ---
the run-to-run variation in $\hat\gamma$ arising because the 5000-cluster dataset
is one random draw from a larger cluster population.

The goal of this study is to verify that a \emph{post-hoc bootstrap} procedure
provides an honest, well-calibrated uncertainty estimate for $\gamma$.  Rather
than re-generating new simulated data, the bootstrap repeats the full analysis on
resampled versions of the \emph{same} fixed dataset.  The study is structured as
two nested levels: $K = 10$ independent outer simulations (each a fresh 5000-cluster
survey), and $B = 10$ bootstrap draws per simulation.  Calibration is assessed by
comparing the scatter of ensemble median estimates across simulations,
$\mathrm{std}_k(\hat\gamma_\mathrm{ens}^{(k)})$, against the average bootstrap
standard deviation, $\mathrm{mean}_k(\sigma_\mathrm{boot}^{(k)})$.  A ratio close
to unity indicates that the bootstrap uncertainty correctly tracks the true
survey-level uncertainty.

% ============================================================
\section{Simulation Pipeline}
% ============================================================

This section describes every step required to reproduce the simulated datasets.

\subsection{Truth catalog}
\label{sec:truth}

A pool of $N_\mathrm{pool} = 5 \times N = 25{,}000$ candidate clusters is generated
as follows.

\begin{enumerate}
  \item \textbf{Redshift.}  Draw $z \sim \mathrm{Uniform}[0.2,\,0.65]$.
        These redshifts are used as an initial sample and subsequently reweighted
        (see below).

  \item \textbf{Halo mass.}  Draw $\log_{10}(M_{200}/M_\odot)$ from the
        Tinker et al.\ (2008) halo mass function in the range
        $\log_{10}(M_{200}/M_\odot) \in [13.5,\,15.0]$.

  \item \textbf{BCG stellar mass.}
        \begin{equation}
          \log_{10}\!\left(\frac{M_\mathrm{BCG}}{M_\odot}\right)
          = 0.4\,\log_{10}\!\left(\frac{M_\mathrm{halo}}{M_\odot}\right)
          + 6
          + \mathcal{N}(0,\,0.15^2).
        \end{equation}

  \item \textbf{Satellite galaxies.}  Drawn from the Zheng et al.\ (2007) halo
        occupation distribution (HOD) calibrated to the SDSS 18 sample.
        Clusters with zero detected satellites are rejected.

  \item \textbf{Magnitude gap.}
        \begin{equation}
          \Delta m = 2.5\,\log_{10}\!\left(\frac{L_\mathrm{BCG}}{L_\mathrm{brightest\;sat}}\right)
          \quad[\mathrm{mag}],
        \end{equation}
        positive by construction.

  \item \textbf{NFW concentration.}  The concentration is modelled as
        \begin{equation}
          \log_{10} c = \log_{10} c_\mathrm{DM14}(z,M)
          + \gamma_\mathrm{true}\,\log_{10}\!\left(\frac{\Delta m}{\Delta m_\mathrm{ref}}\right)
          + \mathcal{N}(0,\,\sigma_{\log c}^2),
        \end{equation}
        with $\sigma_{\log c} = 0$ (no intrinsic scatter in this run),
        $\gamma_\mathrm{true} = 0.2$, and $\Delta m_\mathrm{ref} = 0.1$~mag.
        The Dutton \& Maccio (2014) baseline concentration is
        \begin{equation}
          \log_{10} c_\mathrm{DM14} = 0.905
          - 0.101\,\log_{10}\!\left(\frac{M_{200}\,h}{10^{12}\,M_\odot}\right).
        \end{equation}

  \item \textbf{Resampling to DES Y3 $n(z)$.}  From the 25{,}000-cluster pool,
        draw $N = 5000$ clusters \emph{without replacement}, weighted by
        $P(z) \propto z^2$ on $z \in [0.2,\,0.65]$.  This is a volume-weighted
        approximation to the DES Y3 cluster redshift distribution.
\end{enumerate}

\begin{figure}[htbp]
  \centering
  \includegraphics[width=\linewidth]{../gamma_repeat_resampled_outputs/run_2026-05-23T16-46-18/des_y3_nz.png}
  \caption{DES Y3 cluster and source redshift distributions used in the simulation.}
  \label{fig:des_y3_nz}
\end{figure}

\begin{figure}[htbp]
  \centering
  \includegraphics[width=\linewidth]{../gamma_repeat_resampled_outputs/run_2026-05-23T16-46-18/simulation_000/truth_distributions.png}
  \caption{Truth catalog properties for simulation 000: redshift, log halo mass,
  concentration, and $\Delta m$ distributions.}
  \label{fig:truth_distributions}
\end{figure}

\begin{figure}[htbp]
  \centering
  \includegraphics[width=\linewidth]{../gamma_repeat_resampled_outputs/run_2026-05-23T16-46-18/simulation_000/mc_scatter.png}
  \caption{Mass--concentration scatter coloured by $\Delta m$ (left) and
  concentration vs.\ $\Delta m$ (right) for simulation 000.
  The applied $\gamma = 0.2$ slope is visible as a positive trend in the right panel.}
  \label{fig:mc_scatter}
\end{figure}

\subsection{Triaxial shapes}
\label{sec:triaxial}

Each cluster is assigned random triaxial axis ratios and a random line-of-sight
(LOS) orientation, drawn independently:

\begin{itemize}
  \item Intermediate-to-major axis ratio: $b/a \sim \mathrm{Uniform}[0.5,\,1.0]$.
  \item Minor-to-major axis ratio: $c/a \sim \mathrm{Uniform}[0.3,\,b/a]$
        (ensures $c/a \leq b/a$).
  \item LOS polar angle: $\theta_\mathrm{los} = \arccos(U)$ with
        $U \sim \mathrm{Uniform}[0,\,1]$ (isotropic).
  \item LOS azimuthal angle: $\phi_\mathrm{los} \sim \mathrm{Uniform}[0,\,2\pi]$.
\end{itemize}

The virial semi-major axis $a = r_\mathrm{vir}(M, z)$ is computed using the
\texttt{cluster\_toolkit} package.

\subsection{Miscentering}
\label{sec:miscenter}

Cluster centres may be offset from the true halo centre:

\begin{itemize}
  \item Centred fraction: $F_\mathrm{cen} = 0.64$; miscentered fraction:
        $F_\mathrm{mis} = 1 - F_\mathrm{cen} = 0.36$.
  \item Offset magnitude for miscentered clusters:
        $r_\mathrm{off} \sim \mathrm{Rayleigh}(\sigma_\mathrm{off})$ with
        $\sigma_\mathrm{off} = 0.32\,h^{-1}$~Mpc.
        Using the Planck 2018 value $h = 0.6736$ gives
        $\sigma_\mathrm{off} \approx 0.475$~Mpc.
  \item Offset direction: $\phi_\mathrm{mis} \sim \mathrm{Uniform}[0,\,2\pi]$.
  \item Centred clusters: $r_\mathrm{off} = 0$.
\end{itemize}

\begin{figure}[htbp]
  \centering
  \includegraphics[width=\linewidth]{../gamma_repeat_resampled_outputs/run_2026-05-23T16-46-18/simulation_000/miscentering_offsets.png}
  \caption{Miscentering model: offset magnitude distribution (left) and 2D offset
  positions (right).  36\% of clusters are miscentered.}
  \label{fig:miscentering}
\end{figure}

\subsection{Source catalog and lensing simulation}
\label{sec:sources}

For each cluster $i$ the following procedure is used to generate source galaxies
and measure the $\Delta\Sigma$ profile.

\begin{enumerate}
  \item \textbf{Source positions.}  Galaxy positions are drawn from a
        Poisson process on a square field of width 0.006~rad, at a
        surface density of 30~gal\,arcmin$^{-2}$.

  \item \textbf{Source redshifts.}  Drawn from the Smail distribution
        \begin{equation}
          n(z) \propto \left(\frac{z}{z_0}\right)^\alpha
          \exp\!\left[-\left(\frac{z}{z_0}\right)^\beta\right],
          \quad z \in [0,\,2.5],
        \end{equation}
        with $\alpha = 2$, $\beta = 1.5$, $z_0 = 0.4$.

  \item \textbf{Background cut.}  Sources are retained only if
        $z_s > 1.1\,z_l + 0.2$.

  \item \textbf{Intrinsic ellipticities.}  Drawn from a \texttt{GPriorBA}
        distribution with shape-noise parameter $\sigma = 0.27$.

  \item \textbf{Metacalibration responsivity.}  Per-source responsivity
        $\mathcal{R} \sim \mathrm{LogNormal}(\ln 0.5,\; 0.5)$, clipped to
        $[0.05,\,0.99]$.

  \item \textbf{Triaxial NFW shear.}  The function
        \texttt{apply\_triaxial\_NFW\_shear\_region} from \texttt{xlensing} is
        used to apply the gravitational shear field of the triaxial NFW halo to
        each source.  The cluster centre is placed at position
        $(r_\mathrm{off}\cos\phi_\mathrm{mis},\;r_\mathrm{off}\sin\phi_\mathrm{mis})$
        relative to the field centre.

  \item \textbf{$\Delta\Sigma$ measurement.}  The profile is measured via
        \texttt{metacal\_cluster\_lensing} and the \texttt{signal()} method from
        \texttt{xlensing}.

  \item \textbf{Rejection.}  Clusters with fewer than 20 background sources
        after the redshift cut are discarded.
\end{enumerate}

\begin{figure}[htbp]
  \centering
  \includegraphics[width=\linewidth]{../gamma_repeat_resampled_outputs/run_2026-05-23T16-46-18/simulation_000/sample_profiles.png}
  \caption{Six example $\Delta\Sigma$ profiles.  Triangles indicate miscentered
  clusters.}
  \label{fig:sample_profiles}
\end{figure}

\begin{figure}[htbp]
  \centering
  \includegraphics[width=\linewidth]{../gamma_repeat_resampled_outputs/run_2026-05-23T16-46-18/simulation_000/triaxiality_diagnostic.png}
  \caption{Fractional residuals $(\mathrm{data} - \mathrm{truth\;NFW})/(\mathrm{truth\;NFW})$
  showing the scatter injected by triaxiality.  The vertical dashed line marks
  $R_\mathrm{min} = 0.5$~Mpc.}
  \label{fig:triaxiality}
\end{figure}

\subsection{Radial binning}
\label{sec:bins}

Twelve radial bins are log-spaced between $R = 0.158$~Mpc and $R = 6.31$~Mpc.
Bin centres and edges are computed as follows:

\begin{verbatim}
radii = np.logspace(-0.8, 0.8, 12)        # 12 bin centres [Mpc]
dlog  = (log10(radii[-1]) - log10(radii[0])) / 11
edges = np.logspace(log10(radii[0])  - dlog/2,
                    log10(radii[-1]) + dlog/2, 13)
\end{verbatim}

The per-cluster diagonal covariance for bin $b$ is
\begin{equation}
  \sigma_b^2 =
  \begin{cases}
    \displaystyle\left(\frac{\sigma_\mathrm{shape}\,\bar\Sigma_\mathrm{crit}}{\sqrt{n_b}}\right)^{\!2}
    & n_b \geq 5,\\[8pt]
    10^4\;(M_\odot/\mathrm{pc}^2)^2 & n_b < 5,
  \end{cases}
\end{equation}
where $\sigma_\mathrm{shape} = 0.27$, $\bar\Sigma_\mathrm{crit}$ is the
mean critical surface density for the cluster, and $n_b$ is the number of
background sources in bin $b$.

\textbf{Radial cut for fitting}: only bins with $R \geq 0.5$~Mpc are retained,
leaving 7 of the 12 bins.

\subsection{Cluster observable table}
\label{sec:obs}

In addition to the $\Delta\Sigma$ data vectors and covariance matrices, each
cluster contributes the following per-cluster observables to the likelihood:

\begin{itemize}
  \item \textbf{Richness $\lambda$}: derived by inverting the
        McClintock et al.\ (2018) mass--richness relation.
  \item \textbf{Magnitude gap $\Delta m$}: taken directly from the truth catalog.
  \item \textbf{Population miscentering parameters}: $F_\mathrm{mis} = 0.36$ and
        $\sigma_\mathrm{off} = 0.475$~Mpc, constant for all clusters.
        These are not resampled during the bootstrap.
\end{itemize}

% ============================================================
\section{Analysis Pipeline}
% ============================================================

\subsection{Scaling relation model}
\label{sec:model}

The NFW concentration for cluster $i$ is modelled as
\begin{equation}
  \log_{10} c_i = \log_{10} c_0
  + \alpha\,\bigl(\log_{10} M_i - \log_{10} M_\mathrm{piv}\bigr)
  + \beta\,\log_{10}\!\left(\frac{1+z_i}{1+z_\mathrm{ref}}\right)
  + \gamma\,\log_{10}\!\left(\frac{\Delta m_i}{\Delta m_\mathrm{ref}}\right)
  + \log\_M\_norm,
  \label{eq:scaling}
\end{equation}
with pivot mass $M_\mathrm{piv} = 10^{14.344}\,M_\odot$ (McClintock et al.\ 2018),
$z_\mathrm{ref} = 0.35$, and $\Delta m_\mathrm{ref} = 0.1$~mag.
The predicted $\Delta\Sigma$ profile is computed via \texttt{xlensing}
\texttt{NFW\_Delta\_Sigma}.

Miscentering is modelled as a population mixture:
\begin{equation}
  \Delta\Sigma_{\mathrm{model},i}(r;\theta)
  = (1 - F_\mathrm{mis})\,\Delta\Sigma_\mathrm{NFW}(r;\,c_i,M_i,z_i)
  + F_\mathrm{mis}\,\Delta\Sigma_\mathrm{off}(r;\,c_i,M_i,z_i,\sigma_\mathrm{off}),
  \label{eq:miscenter_model}
\end{equation}
where $\Delta\Sigma_\mathrm{off}$ is the azimuthally-averaged NFW profile
displaced by a Rayleigh-distributed offset radius.

\subsection{Log-likelihood}
\label{sec:likelihood}

For cluster $i$ the residual is
$\mathbf{r}_i(\theta) = \mathbf{d}_i - \mathbf{m}_i(\theta)$,
where $\mathbf{d}_i$ is the 7-element $\Delta\Sigma$ data vector (bins with
$R \geq 0.5$~Mpc) and $\mathbf{m}_i(\theta)$ is the corresponding model
prediction at parameter vector
$\theta = (c_0,\,\alpha,\,\beta,\,\gamma,\,\log\_M\_norm)$.
The covariance $\mathbf{C}_i$ is diagonal (shape noise only).

The joint log-likelihood is
\begin{equation}
  \ell(\theta) = \sum_{i=1}^{N} \ell_i(\theta),
  \qquad
  \ell_i(\theta) = -\tfrac{1}{2}\Bigl[
    k\ln(2\pi)
    + \ln|\mathbf{C}_i|
    + \mathbf{r}_i^\top \mathbf{C}_i^{-1} \mathbf{r}_i
  \Bigr],
  \label{eq:loglike}
\end{equation}
with $k = 7$ radial bins per cluster.  The diagonal structure of $\mathbf{C}_i$
means that $\ln|\mathbf{C}_i| = \sum_b \ln \sigma_{i,b}^2$ and
$\mathbf{r}_i^\top\mathbf{C}_i^{-1}\mathbf{r}_i = \sum_b r_{i,b}^2/\sigma_{i,b}^2$;
the Cholesky decomposition is used in the implementation for numerical stability.

\subsection{Priors}
\label{sec:priors}

\begin{table}[htbp]
  \centering
  \caption{Parameter priors used in the Nautilus sampling.}
  \label{tab:priors}
  \begin{tabular}{lll}
    \toprule
    Parameter & Prior & Motivation \\
    \midrule
    $c_0$ & $\mathcal{N}(4.57,\;0.05^2)$ & DM14 pivot at survey mean \\
    $\alpha$ & $\mathcal{N}(-0.10,\;0.05^2)$ & Mass slope \\
    $\beta$ & $\mathcal{N}(-0.40,\;0.20^2)$ & Redshift slope \\
    $\gamma$ & $\mathcal{N}(0.0,\;3.5^2)$ & Magnitude-gap slope (weak) \\
    $\log\_M\_norm$ & $\mathcal{N}(0.0,\;0.10^2)$ & Mass normalisation \\
    \bottomrule
  \end{tabular}
\end{table}

\subsection{Sampler}
\label{sec:sampler}

The posterior is explored with the Nautilus nested sampler
(Lange 2023).  Settings:
\begin{itemize}
  \item Number of live points: $n_\mathrm{live} = 400$.
  \item Wall-time timeout: 600~s per run.
  \item Rank-1 covariance inflation: $\varepsilon = 0$ (disabled).
  \item Parallelism: pool of 15 workers for the likelihood evaluation.
\end{itemize}
The posterior is extracted via \texttt{sampler.posterior()};
the marginal median and 68\% credible interval for $\gamma$ are computed
from the weighted samples using weighted quantiles.

% ============================================================
\section{Bootstrap Calibration Design}
% ============================================================

\subsection{Motivation}
\label{sec:motivation}

Individual Nautilus posteriors on $\gamma$ correctly capture
\emph{shape-noise uncertainty} --- the variation in $\hat\gamma$ due to finite
numbers of background source galaxies.  They do not, however, capture
\emph{cluster sample variance} --- the variation arising because any realisation
of a 5000-cluster survey is one random draw from the underlying cluster
population.  This causes the posterior widths $\sigma_\mathrm{post}$ to be
systematically narrower than the true uncertainty.

An earlier calibration study (running the pipeline on many independent
simulated datasets) confirmed this: the ratio of the run-to-run scatter of the
posterior median to the mean posterior half-width was
$R = \sigma_\mathrm{runs} / \langle\sigma_\mathrm{post}\rangle \approx 2.16$
for standard runs with fresh data each time.  Reporting the individual posterior
width as the uncertainty on $\gamma$ would therefore under-estimate the true
uncertainty by more than a factor of two.

The bootstrap addresses this by propagating the uncertainty through the full
analysis on many resampled versions of the \emph{same} fixed dataset, without
requiring new simulated data.

\subsection{Nested design}
\label{sec:design}

The calibration study employs two nested levels of randomisation.

\textbf{Outer loop} ($K = 10$ independent simulations):
For each simulation index $k \in \{0,\ldots,9\}$, generate a fresh truth catalog
and lensing dataset following Sections~\ref{sec:truth}--\ref{sec:obs}.
Each outer simulation is a genuinely independent realisation of a 5000-cluster
DES Y3-like survey.

\textbf{Inner loop} ($B = 10$ bootstrap draws per simulation):
Fix the dataset from simulation $k$.  For each bootstrap draw $b \in \{0,\ldots,B-1\}$:

\begin{enumerate}
  \item Draw resample indices with replacement:
        \begin{equation}
          \mathbf{I}^{(b)} = \mathrm{rng.choice}(N,\;N,\;\mathrm{replace}=\mathrm{True}).
        \end{equation}
        The expected fraction of unique clusters is $1 - e^{-1} \approx 63.2\%$
        (observed $\approx 63.5\%$); the remaining $\approx 36.8\%$ are
        out-of-bag and receive zero weight.

  \item Construct the resampled dataset:
        \begin{align}
          \mathbf{d}^{(b)} &= \mathbf{d}[\mathbf{I}^{(b)}],\\
          \mathbf{C}^{(b)} &= \mathbf{C}[\mathbf{I}^{(b)}],\\
          \mathrm{cluster\_table}^{(b)} &= \mathrm{cluster\_table}[\mathbf{I}^{(b)}].
        \end{align}

  \item \textbf{Population constants are not resampled.}
        $F_\mathrm{mis} = 0.36$ and $\sigma_\mathrm{off} = 0.475$~Mpc are
        passed as \texttt{np.full(N, value)} arrays and are identical for
        every bootstrap draw and every simulation.

  \item Run Nautilus on
        $(\mathbf{d}^{(b)},\;\mathbf{C}^{(b)},\;\mathrm{cluster\_table}^{(b)})$
        and extract the posterior median $\hat\gamma^{(k,b)}$.
\end{enumerate}

\subsection{Ensemble statistics per simulation}
\label{sec:ensemble}

From the $B$ bootstrap draws of simulation $k$, two summary statistics are
computed:
\begin{align}
  \hat\gamma_\mathrm{ens}^{(k)} &= \mathrm{median}_b\bigl(\hat\gamma^{(k,b)}\bigr),
  \label{eq:gamma_ens}\\
  \sigma_\mathrm{boot}^{(k)} &= \mathrm{std}_b\bigl(\hat\gamma^{(k,b)}\bigr)
  \quad [\text{ddof}=1].
  \label{eq:sigma_boot}
\end{align}
The ensemble median $\hat\gamma_\mathrm{ens}^{(k)}$ is a robust point estimate
for simulation $k$, and $\sigma_\mathrm{boot}^{(k)}$ is the bootstrap estimate
of the uncertainty on that point estimate.

\subsection{Calibration check across outer simulations}
\label{sec:calcheck}

The calibration check tests whether $\sigma_\mathrm{boot}$ correctly predicts
the scatter of $\hat\gamma_\mathrm{ens}$ across independent surveys:
\begin{equation}
  \boxed{R_\mathrm{cal}
  = \frac{\mathrm{std}_k\!\left(\hat\gamma_\mathrm{ens}^{(k)}\right)}
         {\mathrm{mean}_k\!\left(\sigma_\mathrm{boot}^{(k)}\right)}
  \approx 1.}
  \label{eq:rcal}
\end{equation}
If $R_\mathrm{cal} \approx 1$, the bootstrap uncertainty $\sigma_\mathrm{boot}$
correctly estimates the survey-to-survey scatter of the ensemble estimator, and
can be reported as the uncertainty on $\gamma$ for a single observed survey.

% ============================================================
\section{Results}
% ============================================================

\subsection{Per-simulation summary table}
\label{sec:table}

Table~\ref{tab:results} summarises the key statistics for each of the $K = 10$
outer simulations.  The bias column gives
$(\hat\gamma_\mathrm{ens}^{(k)} - \gamma_\mathrm{true})/\sigma_\mathrm{boot}^{(k)}$
in units of $\sigma_\mathrm{boot}$.

\begin{table}[htbp]
  \centering
  \caption{Per-simulation bootstrap calibration results.
  $\hat\gamma_\mathrm{ens}$: ensemble median;
  $\sigma_\mathrm{boot}$: bootstrap standard deviation of $\hat\gamma$;
  $\langle\sigma_\mathrm{post}\rangle$: mean individual posterior 68\% half-width;
  Bias: $(\hat\gamma_\mathrm{ens} - 0.2)/\sigma_\mathrm{boot}$.}
  \label{tab:results}
  \begin{tabular}{ccccc}
    \toprule
    Sim & $\hat\gamma_\mathrm{ens}$ & $\sigma_\mathrm{boot}$ &
    $\langle\sigma_\mathrm{post}\rangle$ & Bias ($\sigma_\mathrm{boot}$) \\
    \midrule
    0 & $-0.096$ & $0.528$ & $0.149$ & $-0.56$ \\
    1 & $-0.065$ & $0.345$ & $0.126$ & $-0.77$ \\
    2 & $+0.343$ & $0.370$ & $0.162$ & $+0.39$ \\
    3 & $+0.084$ & $0.248$ & $0.103$ & $-0.47$ \\
    4 & $+0.836$ & $0.561$ & $0.106$ & $+1.14$ \\
    5 & $+0.517$ & $0.567$ & $0.126$ & $+0.56$ \\
    6 & $+0.478$ & $0.195$ & $0.152$ & $+1.42$ \\
    7 & $+0.423$ & $0.483$ & $0.191$ & $+0.46$ \\
    8 & $-0.437$ & $0.250$ & $0.112$ & $-2.55$ \\
    9 & $+0.219$ & $0.504$ & $0.163$ & $+0.04$ \\
    \bottomrule
  \end{tabular}
\end{table}

\subsection{Calibration statistics}
\label{sec:calstats}

The aggregate calibration statistics across all $K = 10$ simulations are:

\begin{align}
  \mathrm{mean}_k\!\left(\hat\gamma_\mathrm{ens}^{(k)}\right)
  &= +0.230
  \quad\bigl[\gamma_\mathrm{true} = 0.200;\;\text{mean bias} = +0.030\bigr],
  \label{eq:mean_gamma}\\[4pt]
  \mathrm{std}_k\!\left(\hat\gamma_\mathrm{ens}^{(k)}\right)
  &= 0.368
  \quad\leftarrow\text{scatter of ensemble medians across simulations},
  \label{eq:std_gamma}\\[4pt]
  \mathrm{mean}_k\!\left(\sigma_\mathrm{boot}^{(k)}\right)
  &= 0.405
  \quad\leftarrow\text{mean bootstrap uncertainty estimate},
  \label{eq:mean_boot}\\[4pt]
  \mathrm{mean}_k\!\left(\langle\sigma_\mathrm{post}\rangle^{(k)}\right)
  &= 0.139
  \quad\leftarrow\text{mean individual posterior half-width}.
  \label{eq:mean_post}
\end{align}

The calibration ratio is
\begin{equation}
  \boxed{R_\mathrm{cal}
  = \frac{\mathrm{std}_k(\hat\gamma_\mathrm{ens}^{(k)})}
         {\mathrm{mean}_k(\sigma_\mathrm{boot}^{(k)})}
  = \frac{0.368}{0.405} = 0.91.}
  \label{eq:rcal_result}
\end{equation}

The bootstrap inflation factor relative to the individual posteriors is
\begin{equation}
  \frac{\mathrm{mean}_k(\sigma_\mathrm{boot}^{(k)})}
       {\mathrm{mean}_k(\langle\sigma_\mathrm{post}\rangle^{(k)})}
  = \frac{0.405}{0.139} = 2.91.
  \label{eq:inflation}
\end{equation}

The mean bias of $+0.030$ is well within one-tenth of a $\sigma_\mathrm{boot}$,
confirming that the ensemble median is an approximately unbiased estimator of
$\gamma_\mathrm{true} = 0.2$ when averaged over many independent surveys.

\subsection{Discussion of multimodal posteriors}
\label{sec:multimodal}

The likelihood surface for $\gamma$ exhibits known spurious modes.  Bootstrap
scatter plots (Figures~\ref{fig:gamma_residuals_000} and~\ref{fig:gamma_residuals_004})
reveal two secondary modes at approximately $\hat\gamma \approx +1.0$ to $+1.3$
and $\hat\gamma \approx -0.5$ to $-0.6$, in addition to the primary mode near
$\gamma_\mathrm{true} = 0.2$.

For simulations 4 and 5, a majority of bootstrap draws land in the upper
spurious mode, producing ensemble medians ($\hat\gamma_\mathrm{ens} = +0.836$
and $+0.517$) substantially above $\gamma_\mathrm{true}$.  Despite this,
$R_\mathrm{cal} = 0.91 \approx 1$ because $\sigma_\mathrm{boot}$ correctly
tracks the width of the bootstrap distribution --- whether it is centred on
the correct mode or not.  The bootstrap is therefore \emph{consistent} in the
sense that it accurately reports the uncertainty of the ensemble estimator, even
when that estimator is biased due to modal confusion.

The overlaid posteriors (Figure~\ref{fig:posteriors_overlay}) make the
multi-modal structure visible: individual bootstrap draws contribute distinct
peaks across the three modes, and the ensemble spread faithfully reflects the
resulting uncertainty.

\subsection{Figures}
\label{sec:figures}

\begin{figure}[htbp]
  \centering
  \includegraphics[width=\linewidth]{../gamma_repeat_resampled_outputs/run_2026-05-23T16-46-18/simulation_000/gamma_residuals.png}
  \caption{Bootstrap scatter for simulation 000.  Error bars are individual
  posterior 68\% credible intervals.  The horizontal band shows $\pm\sigma_\mathrm{boot}$.}
  \label{fig:gamma_residuals_000}
\end{figure}

\begin{figure}[htbp]
  \centering
  \includegraphics[width=\linewidth]{../gamma_repeat_resampled_outputs/run_2026-05-23T16-46-18/simulation_004/gamma_residuals.png}
  \caption{Bootstrap scatter for simulation 004 --- an example where most draws
  land in the upper spurious mode near $\hat\gamma \approx +1.1$.}
  \label{fig:gamma_residuals_004}
\end{figure}

\begin{figure}[htbp]
  \centering
  \includegraphics[width=\linewidth]{../gamma_repeat_resampled_outputs/run_2026-05-23T16-46-18/simulation_000/posteriors_overlay.png}
  \caption{Overlaid $\gamma$ posteriors from all 10 bootstrap draws of simulation 000.
  Multiple modes are visible.}
  \label{fig:posteriors_overlay}
\end{figure}

\begin{figure}[htbp]
  \centering
  \includegraphics[width=\linewidth]{../gamma_repeat_resampled_outputs/run_2026-05-23T16-46-18/calibration_summary.png}
  \caption{Main calibration result: ensemble $\hat\gamma$ per simulation (points)
  with $\sigma_\mathrm{boot}$ error bars (lines).  The shaded band is
  $\pm\mathrm{std}(\hat\gamma_\mathrm{ens})$.  The inset reports the calibration ratio.}
  \label{fig:calibration_summary}
\end{figure}

% ============================================================
\section{Conclusions}
% ============================================================

The main findings of this bootstrap calibration study are as follows.

\textbf{The bootstrap is well-calibrated.}
The calibration ratio $R_\mathrm{cal} = 0.91$ is consistent with unity within
the statistical uncertainty of $K = 10$ simulations.  Using the approximation
$\sigma(R_\mathrm{cal}) \approx R_\mathrm{cal}/\sqrt{2(K-1)} \approx 0.21$, the
result $R_\mathrm{cal} = 0.91$ lies within $0.4\,\sigma$ of 1.  The bootstrap
therefore correctly estimates the survey-level uncertainty on $\gamma$, and
$\sigma_\mathrm{boot}$ should be adopted as the reported uncertainty when
analysing a single observed survey.

\textbf{Bootstrap inflation of $\sim 3\times$ over individual posteriors.}
The ratio $\sigma_\mathrm{boot}/\langle\sigma_\mathrm{post}\rangle = 2.91$
confirms that the dominant missing variance in individual Nautilus posteriors is
not measurement (shape) noise but cluster-sample variance and modal variance.
Reporting the individual posterior width as the uncertainty on $\gamma$ would
underestimate the true uncertainty by nearly a factor of three.

\textbf{Multimodality is the main practical limitation.}
The likelihood surface exhibits spurious modes at $\gamma \approx +1.0$--$+1.3$
and $\gamma \approx -0.5$--$-0.6$.  In some survey realisations
(e.g.\ simulations 4 and 5) the majority of bootstrap draws are captured by the
upper spurious mode, producing ensemble medians far from $\gamma_\mathrm{true}$.
Resolving or suppressing these spurious modes --- through tighter priors,
reparametrisation, or a dedicated mode-finding step --- is the highest-priority
next step for the analysis.

\textbf{Bootstrap sample size.}
With $B = 10$ draws the statistical uncertainty on $\sigma_\mathrm{boot}$ itself
is $\sigma(\sigma_\mathrm{boot})/\sigma_\mathrm{boot} \approx 1/\sqrt{2(B-1)} \approx 24\%$.
Increasing to $B = 50$ would reduce this to $\approx 10\%$, providing a much
more stable estimate of $\sigma_\mathrm{boot}$ per simulation.

\end{document}
